\chapter{Sistema de planes y \emph{billing}}\label{cap:billing}

\section{Planes}

\begin{table}[h]
  \centering
  \begin{tabular}{lllll}
    \toprule
    Plan        & Análisis/mes & Verificaciones/mes & Historial & Claves API \\
    \midrule
    \textsc{free}      & limitado & 0  & últimos 30 días & no  \\
    \textsc{pro}       & ampliado & 0  & completo        & 1   \\
    \textsc{ultra}     & alto     & 0  & completo        & 5   \\
    \textsc{super pro} & alto     & sí & completo        & 10  \\
    \bottomrule
  \end{tabular}
  \caption{Características principales por plan.}
  \label{tab:planes}
\end{table}

% TODO: rellenar números exactos de cuotas.

\section{Integración con Stripe}

\subsection{Productos y precios}

Cada plan se modela como un \emph{Product} de Stripe con uno o varios
\emph{Prices} (mensual / anual). Los identificadores de los precios se
inyectan vía variables de entorno
(\texttt{STRIPE\_PRICE\_PRO}, \texttt{STRIPE\_PRICE\_ULTRA}, \dots).

\subsection{Flujo de suscripción}

\begin{enumerate}[leftmargin=*,itemsep=0.2em]
  \item El usuario pulsa \emph{Suscribirse} en el dashboard.
  \item El frontend llama a \texttt{POST /billing/checkout-session}.
  \item El backend crea la sesión Stripe Checkout y devuelve la URL.
  \item Tras el pago, Stripe redirige y emite un \emph{webhook}
        \texttt{checkout.session.completed}.
  \item El backend valida la firma, actualiza la tabla
        \texttt{billing.subscriptions} y registra el periodo.
\end{enumerate}

\subsection{Cancelación y \emph{downgrade}}

Stripe gestiona el ciclo de cancelación. El backend escucha
\texttt{customer.subscription.updated} y
\texttt{customer.subscription.deleted} para mantener la tabla local
sincronizada.

\section{Cuotas y consumo}

Tabla \texttt{usage(user\_id, action, ts)}. Para cada acción tarificada
(\texttt{predict}, \texttt{verify}) se ejecuta una transacción
\texttt{INSERT} dentro de la dependencia FastAPI; un trigger SQL
garantiza que no se rebase la cuota del periodo.

\section{Migraciones}

\path{migrations/001_billing.sql} crea las tablas y políticas RLS. Las
migraciones se aplican manualmente vía \texttt{psql} o, en producción,
desde la consola de Supabase.
